% Vyžaduje balíček: graphicx. Obrázky: analyza/analyticka_cast/obrazky/
% (\graphicspath{{analyza/analyticka_cast/obrazky/}}).
\section{Technická vybavenost}

Technickou vybavenost popisuje v datech pět atributů: připojení na vodovod,
kanalizaci a plyn, způsob vytápění a vybavení výtahem. Všechny jsou vyplněné
společně: alespoň jeden z nich má 68,1\,\% aktuálních objektů a všech pět
64,9\,\%. Vyplněnost se výrazně liší podle způsobu využití
(obrázek~\ref{fig:vybaveni-vyplnenost}). Rodinné a bytové domy je mají
vyplněné ze 100\,\%, občanská vybavenost z 68\,\%, ale garáže jen z 19\,\%
a rodinná rekreace z 23\,\%. Údaje o vybavení tedy RÚIAN eviduje hlavně
u staveb určených k bydlení a ostatní stavby je často nemají; v následujících
grafech jde vždy o podíly z objektů s vyplněnou hodnotou.

\begin{figure}[H]
\centering
\includegraphics[width=0.9\textwidth]{13_vybaveni_vyplnenost_podle_vyuziti}
\caption{Podíl objektů s vyplněnými údaji o technické vybavenosti podle způsobu využití}
\label{fig:vybaveni-vyplnenost}
\end{figure}

\subsection{Přípojky na sítě}

Většina objektů je připojena na veřejné sítě (obrázek~\ref{fig:pripojky}).
Vodovod má 86,0\,\% objektů, 8,8\,\% je bez vodovodu a u 5,2\,\% je hodnota
nezjištěna. Odpadní vody odvádí do kanalizační sítě 78,5\,\% objektů, 9,5\,\%
využívá žumpu, jímku nebo septik, 7,0\,\% nemá kanalizaci ani jímku
a vlastní čističku má jen 0,4\,\%. Plyn z veřejné sítě má 76,7\,\% objektů
a 17,9\,\% je bez plynu; zásobníkový a jiný plyn se vyskytuje zcela
výjimečně (0,1\,\%).

\begin{figure}[H]
\centering
\includegraphics[width=0.9\textwidth]{14_pripojky_voda_kanalizace_plyn}
\caption{Připojení objektů na vodovod, kanalizaci a plyn}
\label{fig:pripojky}
\end{figure}

Brno je lépe vybaveno než okolní obce (obrázek~\ref{fig:vybaveni-brno}).
Rozdíl je malý u vodovodu (92\,\% oproti 87\,\%), ale výrazný
u kanalizační sítě (87\,\% oproti 71\,\%) a plynu z veřejné sítě (85\,\%
oproti 72\,\%). Podíly jsou počítány z objektů se zjištěným údajem.

\begin{figure}[H]
\centering
\includegraphics[width=0.9\textwidth]{15_vybaveni_brno_okoli}
\caption{Podíl objektů s přípojkou nebo výtahem v Brně a okolních obcích}
\label{fig:vybaveni-brno}
\end{figure}

\subsection{Způsob vytápění}

Číselník způsobu vytápění má 39 hodnot. Základní rozlišení je na vytápění
centrální domovní (kotel ve stavbě), centrální dálkové (kotel mimo stavbu)
a lokální; od roku 2024 číselník navíc určuje zdroj energie. Nejčastější je
centrální domovní vytápění (62,5\,\%), následuje lokální (15,5\,\%), jiné
včetně objektů bez topení (11,5\,\%) a centrální dálkové (5,5\,\%).
U 4,9\,\% objektů je druh nezjištěn (obrázek~\ref{fig:vytapeni}).

Zdroj energie je ale evidován jen u 5\,734 objektů, protože starší záznamy
používají původní hodnoty bez zdroje. U těchto objektů převažuje zemní plyn
(54\,\%), následuje elektřina (24\,\%) a tepelné čerpadlo (14\,\%). Jelikož
jde o nově evidované objekty, tento údaj nelze přenášet na celou zástavbu.

\begin{figure}[H]
\centering
\includegraphics[width=\textwidth]{16_zpusob_vytapeni}
\caption{Druh vytápění a zdroj energie (pouze objekty s vyplněným zdrojem)}
\label{fig:vytapeni}
\end{figure}

\subsection{Výtah}

Výtahem je vybaveno jen 3,8\,\% aktuálních objektů; mezi objekty se
zjištěným údajem je to 7,8\,\% v Brně a 1,0\,\% v okolních obcích. Podíl
silně závisí na podlažnosti (obrázek~\ref{fig:vytah}): u objektů do tří
podlaží má výtah nejvýše 2\,\%, u čtyřpodlažních 9\,\%, u pětipodlažních
20\,\%, u šesti- až sedmipodlažních 69\,\% a u objektů s osmi a více
podlažími 97\,\%.

\begin{figure}[H]
\centering
\includegraphics[width=0.9\textwidth]{17_vytah_podle_podlazi}
\caption{Podíl objektů s výtahem podle počtu podlaží}
\label{fig:vytah}
\end{figure}
